\SourceLeaf{482}
\Needspace{6\baselineskip}\begin{center}\bfseries CLOSING NOTE ON THE NEW EDITION.\end{center}

This first volume of \textsc{Gauss}'s works is a reprint of the seven sections of the Disquis. Arithm. published in octavo in 1801. The eighth section, referred to at several places and probably originally intended by \textsc{Gauss} for publication with the others, is among his manuscripts. Since he did not finish preparing it in the same way as the first seven sections, it will accompany the posthumous arithmetical papers in this edition.

Textual changes have been made in the Disqu. Ar. only at the following places:

In Article 125 the two insertions (if \(>5\)) and (\(>17\); but \(-13N3,\allowbreak\ -17N5\)) have been added.

In Article 126. \textit{Proof.} ‘in the series (I), \(a\) terms are divisible by \(p\), \(b\) terms by \(p^2\), \(c\) terms by \(p^3\), etc.’ has been substituted for ‘in the series (I), \(a\) terms are divisible by \(p\) but not by \(p^2\), \(b\) terms by \(p^2\) but not by \(p^3\), \(c\) terms by \(p^3\) but not by \(p^4\), etc.’

In Article 128 III the passage ‘For if \(+b\) or \(-b\equiv r\ (\mathrm{mod.}\ p)\), then \(bb\equiv rr\ (\mathrm{mod.}\ 2p)\), and hence the term \(\frac12(a-bb)\) is divisible by \(p\).’ was still followed, in the edition of 1801, by ‘and much more so \(2(a-bb)\).’; these latter words have been omitted.

In Article 129, \(2\surd a+1\) has everywhere been substituted for \(2\surd a\), and accordingly 9 for 4 as the number below which \(a\) certainly does not fall.

In Article 139 the end of the investigation of case (4) read ‘It is easy to see that from this equation we can deduce \(a'pRh,\allowbreak\ \pm ahpRa',\allowbreak\ \pm aa'hRp\); these agree with what we found in (2). In the remaining cases the proof is the same.’ This has been changed in accordance with the note that \textsc{Gauss} added to Article 2 of the paper \textit{Theorematis arithmetici demonstratio nova}: ‘It will not be out of place to indicate and correct here a slight error that, through I know not what negligence, slipped into the exposition of that proof. On page (108), from line (14) onward, the following reasoning is to be substituted: It is easy to see that from this equation we can deduce \(a'pRh\ldots(\alpha);\allowbreak\ \pm ahRa'\ldots(\delta);\allowbreak\ \pm ahRp\ldots(\gamma)\). From \((\alpha)\), which agrees with \((\alpha)\) in (2), it follows, just as there, that either \(hRp,\allowbreak\ hRa'\) hold together, or \(hNp,\allowbreak\ hNa'\). But in the former case \((\delta)\) would give \(aRa'\), contrary to the hypothesis; hence \(hNp\), and therefore by \((\gamma)\) also \(aNp\).

\SourceLeaf{483}
In Article 171, in accordance with a manuscript note, ‘in which \(A\) is neither greater than \(\surd{\frac43D},\allowbreak\ C\), nor less than \(2B\)’ has been substituted for ‘in which \(A\) is not \(>\surd{\frac43D}\), \(B\) not \(>\frac12A\), \(C\) not \(<A\)’.

In Article 174. \textit{Second method.} ‘Since \(a\) will not be \(>\surd{\frac43D}\), all forms obtained in this way will evidently be reduced.’ has been substituted for ‘Any forms obtained in this way in which \(a>\surd{\frac43D}\) must be rejected; all the rest will evidently be reduced.’

In Article 256 VI the earlier edition includes, among the positive determinants of the form \(8n+5\) for which the number of classes in the properly primitive order is three times as large as that in the improperly primitive order, also the determinant 397; but it has equally many classes of both kinds.

In Article 302, in accordance with a manuscript note, the number of classes for the second thousand of negative determinants is given as 28595, rather than 28603 as in the earlier edition.

\Needspace{6\baselineskip}In Article 303 the numbers of determinants for the classifications II. 4; II. 5; IV. 2 are given, from the tables found among the posthumous papers, as 31, 44, 69, rather than 32, 42, 68 as in the earlier edition.

In Article 325, ‘\(23t,\allowbreak\ 23t\pm1,\allowbreak\ 23t\pm5,\allowbreak\ 23t\pm7,\allowbreak\ 23t\pm9,\allowbreak\ 23t\pm10\). When these have been deleted, 119, 127 are found to remain, and these two alone give \(V\) a square value’ has been substituted for ‘\(23t,\allowbreak\ 23t\pm5,\allowbreak\ 23t\pm7\), \(23t\pm9,\allowbreak\ 23t\pm10\). When these have been deleted, 119, 127, 137 are found to remain, of which the first two alone give \(V\) a square value.’

In Article 360 IV the words ‘whereas for most other cubic equations whose roots are all real, trisection of both an angle and a ratio cannot be avoided.’, which followed ‘for example, for \(\beta=3\) only trisection of an angle’, and whose incorrectness \textsc{Gauss} noted in his own copy, have been omitted.

The notes on pages 80, 102, 133, 211, 263, 265, 449 are taken from the posthumous manuscripts.

To make the work easier to survey, the outward form of the new edition has been changed in several respects from the printing of 1801. This was thought all the more justified because \textsc{Gauss}, by his own account, also considered saving space in other respects. Many formulas that were contained in the line of text have been set out separately. The summaries, which preceded the work as a collection, have also been placed before the individual divisions, rather than only before the sections as in the first edition; the more general ones have been used as page headings.

\BookRule{45mm}

\begin{center}\small GÖTTINGEN,\\ PRINTED AT THE DIETERICH UNIVERSITY PRESS\\ W. FR. KAESTNER.\end{center}
